% Part: proof-theory
% Chapter: normalization
% Section: permutations

\documentclass[../../../include/open-logic-section]{subfiles}

\begin{document}

\olfileid{pt}{nor}{per}

\olsection{تحويلات التبديل}

القطوع في~\Log{N1i} هي المقاطع التي تنتهي بالمقدمة الكبرى لقاعدة
!!a{elimination}. وإحدى حالات برهان التطبيع هي تقصير القطوع القصوى، ومن ثم
تقليل طول القطع في~!!{proof}. ويمكن لقطع طوله أكبر من~$1$ أن ينتهي بطرائق
متعددة، بحسب ما إذا كان آخر ظهور لـ~!!{formula}~$!A$ في القطع نتيجةً لـ
(\Elim{\lor} أو~\Elim{\lexists})، وبحسب استدلال~!!{elimination} الذي تكون
تلك الصيغة مقدمته الكبرى.

فمثلًا، إذا كان الاستدلالان المعنيان هما~\Elim{\lor} و\Elim{\land}، وكانت
$!A \ident !B \land !C$، وكان هناك~!!{proof} جزئي~$\delta_1$ ينتهي هكذا:
\[
\AxiomC{}
\RightLabel{$\delta_2$}
\DeduceC{$!D \lor !E$}
\AxiomC{$\Discharge{!D}{x}$}
\RightLabel{$\delta_3$}
\DeduceC{$!B \land !C$}
\AxiomC{$\Discharge{!E}{y}$}
\RightLabel{$\delta_4$}
\DeduceC{$!B \land !C$}
\DischargeRule{\Elim{\lor}}{x\,y}
\TrinaryInfC{$!B \land !C$}
\RightLabel{\Elim{\land}[1]}
\UnaryInfC{$!B$}
\DisplayProof
\]
فإن آخر ظهور لـ~!!{formula}~$!A_n$ في القطع هو نتيجة~\Elim{\lor}، أي
الظهور السفلي لـ~$!B \land !C$. أما الظهور قبل الأخير لـ~!!{formula}
$!A_{n-1}$ فهو إحدى المقدمتين الصغريين لـ~\Elim{\lor}.

ولتقصير هذا القطع، يمكننا عكس ترتيب استدلالي~\Elim{\lor} و\Elim{\land}
(أي ``تبديلهما''). وبعبارة أخرى، نستبدل بالبرهان الجزئي~$\delta_1$ داخل
$\delta$ ما يأتي:
\[
\AxiomC{}
\RightLabel{$\delta_2$}
\DeduceC{$!D \lor !E$}
\AxiomC{$\Discharge{!D}{x}$}
\RightLabel{$\delta_3$}
\DeduceC{$!B \land !C$}
\RightLabel{\Elim{\land}[1]}
\UnaryInfC{$!B$}
\AxiomC{$\Discharge{!E}{y}$}
\RightLabel{$\delta_4$}
\DeduceC{$!B \land !C$}
\RightLabel{\Elim{\land}[1]}
\UnaryInfC{$!B$}
\DischargeRule{\Elim{\lor}}{x\,y}
\TrinaryInfC{$!B$}
\DisplayProof
\]
فالقطوع التي كانت تنتهي في مقدمة~\Elim{\land} تحت~\Elim{\lor} تنتهي الآن
في مقدمة أحد استدلالي~\Elim{\land} الواقعين فوق المقدمتين الصغريين لـ
\Elim{\lor}؛ أي إن طول كل منها قد نقص بمقدار~$1$.

وفوق ذلك، لا يتغير أي قطع يقع كله داخل~$\delta_2$ أو~$\delta_3$ أو
$\delta_4$، أو يقع كله خارج~$\delta_1$ داخل~$\delta$: فهو ما يزال يشتمل
على~!!{formula}s نفسها ويمر بالاستدلالات نفسها، ولذلك له خصوصًا رتبة القطع
وطوله نفسيهما كما في القطع المقابل في~$\delta$. ومن ثم لم تزد رتبة القطع
في~!!{proof} الجديد، وطول القطع في~!!{proof} الجديد أقل منه في~!!{proof}
القديم.

\begin{prob}
ينقص طول قطع واحد على الأقل بمقدار~$1$ بعد تحويل تبديل. وإذا بدّلنا
\Elim{\land} عبر~\Elim{\lor}، فهناك في الواقع مقطعا قطع ينقص طولهما، ومن ثم
ينقص طول القطع في~!!{proof} بمقدار~$2$ على الأقل. فهل يمكن أن ينقص طول
القطع في~!!a{proof} بأكثر من~$2$ بتحويل واحد من هذا النوع؟ وهل يمكن أن
تنقص \emph{رتبة} القطع؟
\end{prob}

تسمى عملية نقل~!!a{elimination} إلى فوق قاعدة~\Elim{\lor} أو
\Elim{\lexists} \emph{تحويل تبديل}. وترد أنماط بعض أزواج القواعد في
\olref{tab:permutation-conversions}. (وتُترك الأنماط الباقية تمارين.)

% \begin{table}
% \begin{multline*}
% \shoveleft{\AxiomC{$!D \lor !E$}
% \AxiomC{$\Discharge{!D}{x}$}
% \shortDeduce
% \DeduceC{$!B \land !C$}
% \AxiomC{$\Discharge{!E}{y}$}
% \shortDeduce
% \DeduceC{$!B \land !C$}
% \DischargeRule{\Elim{\lor}}{x\,y}
% \TrinaryInfC{$!B \land !C$}
% \RightLabel{\Elim{\land}[1]}
% \UnaryInfC{$!B$}
% \DisplayProof}\\
% \shoveright{\longrightarrow\ 
% \AxiomC{$!D \lor !E$}
% \AxiomC{$\Discharge{!D}{x}$}
% \shortDeduce
% \DeduceC{$!B \land !C$}
% \RightLabel{\Elim{\land}[1]}
% \UnaryInfC{$!B$}
% \AxiomC{$\Discharge{!E}{y}$}
% \shortDeduce
% \DeduceC{$!B \land !C$}
% \RightLabel{\Elim{\land}[1]}
% \UnaryInfC{$!B$}
% \DischargeRule{\Elim{\lor}}{x\,y}
% \TrinaryInfC{$!B$}
% \DisplayProof}
% \\
% \shoveleft{\AxiomC{$!D \lor !E$}
% \AxiomC{$\Discharge{!D}{x}$}
% \shortDeduce
% \DeduceC{$!B \lif !C$}
% \AxiomC{$\Discharge{!E}{y}$}
% \shortDeduce
% \DeduceC{$!B \lif !C$}
% \DischargeRule{\Elim{\lor}}{x\,y}
% \TrinaryInfC{$!B \lif !C$}
% \AxiomC{$!B$}
% \RightLabel{\Elim{\lif}}
% \BinaryInfC{$!C$}
% \DisplayProof}\\
% \shoveright{\longrightarrow\ 
% \AxiomC{$!D \lor !E$}
% \AxiomC{$\Discharge{!D}{x}$}
% \shortDeduce
% \DeduceC{$!B \lif !C$}
% \AxiomC{$!B$}
% \RightLabel{\Elim{\lif}}
% \BinaryInfC{$!C$}
% \AxiomC{$\Discharge{!E}{y}$}
% \shortDeduce
% \DeduceC{$!B \lif !C$}
% \AxiomC{$!B$}
% \RightLabel{\Elim{\lif}}
% \BinaryInfC{$!C$}
% \DischargeRule{\Elim{\lor}}{x\,y}
% \TrinaryInfC{$!C$}
% \DisplayProof}
% \\
% \shoveleft{\AxiomC{$!D \lor !E$}
% \AxiomC{$\Discharge{!D}{x}$}
% \shortDeduce
% \DeduceC{$!B \lif !C$}
% \AxiomC{$\Discharge{!E}{y}$}
% \shortDeduce
% \DeduceC{$!B \lif !C$}
% \DischargeRule{\Elim{\lor}}{x\,y}
% \TrinaryInfC{$!B \lif !C$}
% \AxiomC{$!B$}
% \RightLabel{\Elim{\lif}}
% \BinaryInfC{$!C$}
% \DisplayProof}\\
% \shoveright{\longrightarrow\ 
% \AxiomC{$!D \lor !E$}
% \AxiomC{$\Discharge{!D}{x}$}
% \shortDeduce
% \DeduceC{$!B \lif !C$}
% \AxiomC{$!B$}
% \RightLabel{\Elim{\lif}}
% \BinaryInfC{$!C$}
% \AxiomC{$\Discharge{!E}{y}$}
% \shortDeduce
% \DeduceC{$!B \lif !C$}
% \AxiomC{$!B$}
% \RightLabel{\Elim{\lif}}
% \BinaryInfC{$!C$}
% \DischargeRule{\Elim{\lor}}{x\,y}
% \TrinaryInfC{$!C$}
% \DisplayProof}
% \\
% \shoveleft{\AxiomC{$!D \lor !E$}
% \AxiomC{$\Discharge{!D}{x}$}
% \shortDeduce
% \DeduceC{$\lforall[x][!B(x)]$}
% \AxiomC{$\Discharge{!E}{y}$}
% \shortDeduce
% \DeduceC{$\lforall[x][!B(x)]$}
% \DischargeRule{\Elim{\lor}}{x\,y}
% \TrinaryInfC{$\lforall[x][!B(x)]$}
% \RightLabel{\Elim{\lforall}}
% \UnaryInfC{$!B(t)$}
% \DisplayProof}\\
% \shoveright{\longrightarrow\ 
% \AxiomC{$!D \lor !E$}
% \AxiomC{$\Discharge{!D}{x}$}
% \shortDeduce
% \DeduceC{$\lforall[x][!B(x)]$}
% \RightLabel{\Elim{\lforall}}
% \UnaryInfC{$!B(t)$}
% \AxiomC{$\Discharge{!E}{y}$}
% \shortDeduce
% \DeduceC{$\lforall[x][!B(x)]$}
% \RightLabel{\Elim{\lforall}}
% \UnaryInfC{$!B(t)$}
% \DischargeRule{\Elim{\lor}}{x\,y}
% \TrinaryInfC{$!B(t)$}
% \DisplayProof}
% \end{multline*}
% \caption{Some permutation conversions for \Log{N1i}.}
% \ollabel{tab:permutation-conversions}
% \end{table}

{\small
\begin{longtable}{@{}l@{}}
\AxiomC{$!D \lor !E$}
\AxiomC{$\Discharge{!D}{x}$}
\shortDeduce
\DeduceC{$!B \land !C$}
\AxiomC{$\Discharge{!E}{y}$}
\shortDeduce
\DeduceC{$!B \land !C$}
\DischargeRule{\Elim{\lor}}{x\,y}
\TrinaryInfC{$!B \land !C$}
\RightLabel{\Elim{\land}[1]}
\UnaryInfC{$!B$}
\DisplayProof
$\longrightarrow$\\[6ex]
\quad\AxiomC{$!D \lor !E$}
\AxiomC{$\Discharge{!D}{x}$}
\shortDeduce
\DeduceC{$!B \land !C$}
\RightLabel{\Elim{\land}[1]}
\UnaryInfC{$!B$}
\AxiomC{$\Discharge{!E}{y}$}
\shortDeduce
\DeduceC{$!B \land !C$}
\RightLabel{\Elim{\land}[1]}
\UnaryInfC{$!B$}
\DischargeRule{\Elim{\lor}}{x\,y}
\TrinaryInfC{$!B$}
\DisplayProof
\\[10ex]
\AxiomC{$!D \lor !E$}
\AxiomC{$\Discharge{!D}{x}$}
\shortDeduce
\DeduceC{$!B \lif !C$}
\AxiomC{$\Discharge{!E}{y}$}
\shortDeduce
\DeduceC{$!B \lif !C$}
\DischargeRule{\Elim{\lor}}{x\,y}
\TrinaryInfC{$!B \lif !C$}
\AxiomC{$!B$}
\RightLabel{\Elim{\lif}}
\BinaryInfC{$!C$}
\DisplayProof
$\longrightarrow$\\[6ex]
\AxiomC{$!D \lor !E$}
\AxiomC{$\Discharge{!D}{x}$}
\shortDeduce
\DeduceC{$!B \lif !C$}
\AxiomC{$!B$}
\RightLabel{\Elim{\lif}}
\BinaryInfC{$!C$}
\AxiomC{$\Discharge{!E}{y}$}
\shortDeduce
\DeduceC{$!B \lif !C$}
\AxiomC{$!B$}
\RightLabel{\Elim{\lif}}
\BinaryInfC{$!C$}
\DischargeRule{\Elim{\lor}}{x\,y}
\TrinaryInfC{$!C$}
\DisplayProof
\\[10ex]
\AxiomC{$!D \lor !E$}
\AxiomC{$\Discharge{!D}{x}$}
\shortDeduce
\DeduceC{$\lforall[x][!B(x)]$}
\AxiomC{$\Discharge{!E}{y}$}
\shortDeduce
\DeduceC{$\lforall[x][!B(x)]$}
\DischargeRule{\Elim{\lor}}{x\,y}
\TrinaryInfC{$\lforall[x][!B(x)]$}
\RightLabel{\Elim{\lforall}}
\UnaryInfC{$!B(t)$}
\DisplayProof
\quad$\longrightarrow$\\[8ex]
\quad\AxiomC{$!D \lor !E$}
\AxiomC{$\Discharge{!D}{x}$}
\shortDeduce
\DeduceC{$\lforall[x][!B(x)]$}
\RightLabel{\Elim{\lforall}}
\UnaryInfC{$!B(t)$}
\AxiomC{$\Discharge{!E}{y}$}
\shortDeduce
\DeduceC{$\lforall[x][!B(x)]$}
\RightLabel{\Elim{\lforall}}
\UnaryInfC{$!B(t)$}
\DischargeRule{\Elim{\lor}}{x\,y}
\TrinaryInfC{$!B(t)$}
\DisplayProof
\\
\caption{بعض تحويلات التبديل في~\Log{N1i}.}
\ollabel{tab:permutation-conversions}
\end{longtable}
}

\begin{prob}
اكتب تحويلات التبديل لـ~\Elim{\lor} متبوعةً بـ~\Elim{\lor} أو
\Elim{\lnot}.
\end{prob}

\begin{prob}
اكتب تحويلات التبديل لـ~\Elim{\lexists} متبوعةً بـ~\Elim{\lexists}.
واشرح لماذا لا يُنتهك شرط المتغير الذاتي في الحالة الثانية.
\end{prob}

ينبغي أن نتحقق من أن تطبيق تحويل تبديل لا ينتهك شرط متغير ذاتي. انظر في:
\[
\AxiomC{}
\RightLabel{$\delta_2$}
\DeduceC{$!D \lor !E$}
\AxiomC{$\Discharge{!D}{x}$}
\RightLabel{$\delta_3$}
\DeduceC{$\lexists[x][!B(x)]$}
\AxiomC{$\Discharge{!E}{y}$}
\RightLabel{$\delta_4$}
\DeduceC{$\lexists[x][!B(x)]$}
\DischargeRule{\Elim{\lor}}{x\,y}
\TrinaryInfC{$\lexists[x][!B(x)]$}
\AxiomC{$\Discharge{!B(c)}{z}$}
\RightLabel{$\delta_5$}
\DeduceC{$!C$}
\DischargeRule{\Elim{\lexists}}{z}
\BinaryInfC{$!C$}
\DisplayProof
\]
وقبل استبدال هذا البرهان نأخذ نسختين من~$\delta_5$ مع إعادة تسمية
المتغيرات الذاتية ووسوم التفريغ، تجنبًا للأسر، بحيث تكون الأزواج
$(c_D,z_D)$ و$(c_E,z_E)$ حديثةً ومتميزةً زوجيًا. ثم نستبدله بما يأتي:
\[
\AxiomC{}
\RightLabel{$\delta_2$}
\DeduceC{$!D \lor !E$}
\AxiomC{$\Discharge{!D}{x}$}
\RightLabel{$\delta_3$}
\DeduceC{$\lexists[x][!B(x)]$}
\AxiomC{$\Discharge{!B(c_D)}{z_D}$}
\RightLabel{$\delta_5^{(D)}$}
\DeduceC{$!C$}
\DischargeRule{\Elim{\lexists}}{z_D}
\BinaryInfC{$!C$}
\AxiomC{$\Discharge{!E}{y}$}
\RightLabel{$\delta_4$}
\DeduceC{$\lexists[x][!B(x)]$}
\AxiomC{$\Discharge{!B(c_E)}{z_E}$}
\RightLabel{$\delta_5^{(E)}$}
\DeduceC{$!C$}
\DischargeRule{\Elim{\lexists}}{z_E}
\BinaryInfC{$!C$}
\DischargeRule{\Elim{\lor}}{x\,y}
\TrinaryInfC{$!C$}
\DisplayProof
\]
قد نقلق من أن المتغير الذاتي~$c$ يمكن أن يظهر، مثلًا، في~$!D$. فالفرض
$!D$ يفرَّغ فوق~\Elim{\lexists} الأصلي، ولذلك لا يظهر في~!!{proof} الأصلي
ضمن فرض مفتوح فوق المقدمة الكبرى لاستدلال~\Elim{\lexists}، ويتحقق هناك شرط
المتغير الذاتي. أما في~!!{proof} الجديد فلا يفرَّغ~$!D$ إلا بعد استدلالي
\Elim{\lexists}، فيبدو أن الشرط سيُنتهك. لكننا نفترض أن~!!{proof}sنا
منتظمة: فكل متغير ذاتي هو المتغير الذاتي لاستدلال واحد، ولا يظهر في
!!{proof} إلا فوق المقدمة الصغرى لـ~\Elim{\lexists}. وعلى وجه الخصوص لا
يمكن أن يظهر~$c$ في~$\delta_3$ أو~$\delta_4$؛ وإعادة التسمية الحديثة التي
أجريناها تمنع الأسر بين النسختين أيضًا.

ويبيّن هذا المثال أمرًا آخر: ففي~!!{proof} الجديد نسختان من~$\delta_5$.
ولهذا، إذا احتوى~$\delta_5$ على قطع أقصى، فلا يلزم أن ينقص طول القطع؛ بل
قد يزداد بسبب نسخه. فعلينا أن نطبّق تحويل التبديل فقط حين نعلم أنه لا يوجد
قطع أقصى في~$\delta_5$. ونضمن ذلك باختيار قطع أقصى \emph{علوي} دائمًا؛ أي
قطع أقصى لا يحتوي~!!{proof}ه الجزئي (وهو~$\delta_1$ هنا) قطعًا أقصى ينتهي
فوق آخر استدلال في ذلك~!!{proof} الجزئي. ويستبعد هذا وجود قطوع قصوى أخرى
في~$\delta_5$ (بل وفي~$\delta_2$ أو~$\delta_3$ أو~$\delta_4$ كذلك).

\end{document}
